\section{Background and related work}\label{sec:background}

This section fixes the code, the decoder, the prior work we position against, and a first-order model of the round-one window.

\subsection{QC-MDPC codes and the BF decoder}

Quasi-cyclic moderate-density parity-check (QC-MDPC) codes are binary linear codes with block-circulant parity checks \cite{misoczki2013mdpc}. We write $r$ for the block size and $n=2r$ for the code length. The parity-check matrix is $H=[H_0\mid H_1]$, where each circulant block has row weight $d$; the full matrix has row weight $2d$ and column weight $d$. The symbol $t$ is the number of errors injected into one decoding attempt. The primary operating point is $(r,d,t)=(4801,45,92)$ (Section~\ref{sec:method}). For a parity-check matrix $H$ and an error vector $e$, the syndrome is $s=He^{\top}$. A decoding attempt succeeds when the final syndrome is zero and fails otherwise.

BF decoding repairs errors iteratively. The decision variable of a position is its unsatisfied parity-check counter (UPC), the number of unsatisfied equations that involve it. In every round the decoder flips all positions whose UPC passes a threshold, recomputes the syndrome, and repeats until the syndrome vanishes or the iteration budget ends. Algorithm~\ref{alg:bf} states the procedure.

\begin{algorithm}[t]
\caption{Plain BF decoding}
\label{alg:bf}
\begin{algorithmic}[1]
\Require $H$, input syndrome $s_{\mathrm{in}}=He^{\top}$, budget $I$
\Ensure error estimate $\hat e$, or failure
\State $\hat e \gets \mathbf{0}$; $s_{\mathrm{res}} \gets s_{\mathrm{in}}$
\For{$i=1$ to $I$}
  \State compute $U_j=\mathrm{UPC}(H,s_{\mathrm{res}},\hat e,j)$ for every position $j$
  \State set $\tau_i$ from the current syndrome weight $|s_{\mathrm{res}}|$
  \State flip every $j$ with $U_j\ge\tau_i$
  \State update $\hat e$ and $s_{\mathrm{res}}\gets s_{\mathrm{in}}\oplus H\hat e^{\top}$
  \If{$s_{\mathrm{res}}=\mathbf{0}$}
    \State \textbf{return} $\hat e$
  \EndIf
\EndFor
\State \textbf{return} failure
\end{algorithmic}
\end{algorithm}

The baseline is the CLASSIC BF decoder of Vasseur's reference implementation \cite{vasseur2021thesis}, denoted C0. Its threshold is the binomial-model value of \path{compute_threshold}, rounded up and clipped to $[(d+1)/2,d]$, and recomputed from the current syndrome weight in every round. The main experiment uses a 20-round budget and static-distance semantics. Static runs use a fixed initial $t$; the offset is added to the clipped base threshold and the result is clipped again to the same interval. The archived note \path{reference/FROZEN_REFERENCE_SEMANTICS.md} shipped with the code gives the exact formula and source identity. A run stops on zero syndrome or budget exhaustion. A separately coded dynamic-distance reference uses the current distance to the true error, information a deployer does not have, and serves only as a simulation-privileged directional check. Its named descendants modify the schedule or carry additional state: Black-Gray-Flip (BGF), Backflip, and the specification's decoder, which we call BIKE-Flip \cite{drucker2020shades,vasseur2021thesis,bike2024spec}. We treat them as separate decoders. BIKE-Flip protects every round with a floor, $T=\max(f_t(S),T_i(S_0))$ \cite{bike2024spec}, where $S$ is the current syndrome weight, $f_t(S)$ an affine function of it, and $T_i(S_0)$ a floor computed from the initial syndrome weight $S_0$. Each floor equals the un-margined level plus a margin $\Delta$, so the margin persists in every iteration rather than appearing only in the first. At Level 1 the specification sets $\Delta=3$ and a budget of seven iterations. Appendix A.2.4 derives the affine constants as the tangent of the theoretical threshold at the average syndrome weight; we verified our implementation against the constants printed for Levels 1, 3 and 5. The published rationale places the residual failures at iterations two and three, where the falling syndrome weight pushes the threshold below the counters of a few correct positions \cite{sendrier2024seminar,wang2023weak}. FRO-2 takes a different route: it adds $+2$ to the round-one threshold once, and later rounds keep the plain threshold.

\subsection{Failure is conditional on keys and schedules}

Failure is conditional on the key as well as the channel. Guo et al.\ showed that error patterns triggering failure correlate with the secret matrix, so observed failures leak key information \cite{guo2016key}; Wang et al.\ turned the same correlation into a weak-key construction \cite{wang2023weak}; and Eaton et al.\ converted timing observations into a key-recovery attack \cite{eaton2018timing}. Gellersen et al.\ extended the side-channel route \cite{gellersen2025}. Weak-key classes with elevated failure rates have also been constructed, enumerated and filtered \cite{vasseur2021thesis,nosouhi2023weakkey,matthews2025combinatorial}. The schedule is a second conditioning variable, because the flips of one round shape the state entering the next. Our design therefore perturbs the schedule and treats the key as the unit of inference: all variants decode identical $(H,e,s)$ instances, and uncertainty comes from resampling keys rather than errors (Section~\ref{sec:method}).

\subsection{Related work}

Failure modelling has followed two routes. The first explains failure from decoder state. Kim and Park fit a per-instance logistic predictor \cite{kim2021logistic}; Paiva and Terada use the first-round residual error as a proxy \cite{paiva2022pickyfix}; Santini et al.\ and Annechini et al.\ analyse counters for regular codes \cite{santini2019conservative,annechini2024estimation}. The second estimates rates beyond simulation. Sendrier and Vasseur model BF as a Markov chain \cite{sendrier2019dfr,sendrier2020lowdfr}; Tillich derives failure-probability bounds \cite{tillich2018dfp}; Vasseur adds extrapolation with confidence intervals \cite{vasseur2021thesis}; Arpin et al.\ predict key-dependent error floors \cite{arpin2025errorfloor}; and Annechini et al.\ give a closed form with an explicit warning that proxy-based extrapolation is dangerous \cite{annechini2025closedform}.

A third line tunes thresholds directly. Kichna et al.\ use particle-swarm search at our block size, reporting mean iterations reduced from 5.36 to 3.01 at the recommended error weight, but no failure-rate measurement \cite{kichna2025pso}. Recent work changes the decoder itself: BF-Max \cite{baldelli2025bfmax}, in-place randomised flipping \cite{baldi2021inplace}, near-codeword-aware BF \cite{baldelli2026ncw}, and follow-up refinements \cite{santini2026flip}. These decoders are reference points for magnitude and cost, not objects we claim to dominate. Taken together, these lines motivate the question this paper pursues: does a small change in the round-one schedule alter conditional failure, and does the effect depend on the round and persistence? We test that question on paired instances; Section~\ref{sec:counterfactual} adds a secondary diagnostic, asking whether the same first-round statistics can predict another schedule.
